%=============================================================================
% AB INITIO CODE ANALYSIS
% What computation does the Python code ACTUALLY perform?
%=============================================================================
\documentclass[11pt]{article}
\usepackage{amsmath,amssymb,booktabs,geometry,listings,xcolor}
\geometry{margin=1in}
\lstset{language=Python, basicstyle=\ttfamily\small,
  keywordstyle=\color{blue}, commentstyle=\color{gray},
  breaklines=true, frame=single}

\title{Ab Initio Code Analysis:\\
  What Computation Produces $\alpha^{-1} = 137.03599985$?}
\date{22 March 2026}

\begin{document}
\maketitle

%=============================================================================
\section{Summary of Findings}
%=============================================================================

\textbf{Answer: (c) Direct algebraic formula, but with one critical gap.}

The published code in the Ab Initio paper contains TWO code listings.
Neither performs a lattice Monte Carlo simulation or a spectral action /
heat kernel computation. The code that produces $\alpha^{-1} = 137.03599985$
uses a \textbf{direct algebraic formula} that combines:
\begin{enumerate}
\item A finite weighted sum (the Chern--Simons average $\beta_{U(1)}$),
\item Fixed rational correction factors $(63/64)$ and $(4096/4095)$,
\item The electroweak formula $\alpha^{-1} = 6\pi\,\kappa_{U(1)}$.
\end{enumerate}

However, the \textbf{exact sub-ppm value} $\kappa_{U(1)} = 7.26999\ldots$
is \emph{stated as a result} of a spectral action computation on
$\mathbb{C}P^2 \times S^3$---but that spectral action computation is
\textbf{never shown} in any published code or any tex file in the corpus.
The code that exists either reverse-engineers $\kappa_{U(1)}$ from the
target $\alpha^{-1}$, or uses an approximate closed-form formula accurate
to $\sim 85$~ppm.

%=============================================================================
\section{The Two Code Listings in the Ab Initio Paper (Pages 23--24)}
%=============================================================================

\subsection{Listing 1: UV Cutoff $k_{\max} = 60$}

This code computes $\beta_{U(1)} = \langle k+2 \rangle$ for various
$k_{\max}$ values. The exact code is:

\begin{lstlisting}
import math

def w(k):
    """Microsector weight from SU(2) CS on S^3 (Witten 1989)"""
    return (2.0/(k+2)) * (math.sin(math.pi/(k+2)))**2

for k_max in [50, 60, 100, 1000000]:
    Z = sum(w(k) for k in range(k_max))
    k_eff = sum((k+2)*w(k) for k in range(k_max)) / Z
    print(f"k_max={k_max:7d}: <k+2> = {k_eff:.4f}")
\end{lstlisting}

\textbf{Output} (as printed in the paper):
\begin{verbatim}
k_max=     50: <k+2> = 3.7705
k_max=     60: <k+2> = 3.7969   <-- Derived from Bridge Lemma
k_max=    100: <k+2> = 3.8517
k_max=1000000: <k+2> = 3.9386   <-- converged; gives alpha = 1/303,
                                      ruled out
\end{verbatim}

This listing demonstrates that:
\begin{itemize}
\item At $k_{\max} = 60$: $\beta_{U(1)} = 3.7969\ldots$ (precise value: $3.796945276\ldots$)
\item The converged infinite sum gives $\beta \approx 3.94$, yielding $\alpha = 1/303$
\item Only the truncated sum at $k_{\max} = 60$ gives the correct physical value
\end{itemize}

\subsection{Listing 2: Wilson-Normalized $\alpha$ from Stiffnesses}

\begin{lstlisting}
def alpha_wilson(ku, ks):
    g1 = 1.0/ku           # U(1): standard
    g2 = 4.0/ks           # SU(2): Wilson normalization (beta = 4/g^2)
    e2 = g1*g2/(g1+g2)
    return e2/(4.0*math.pi)
\end{lstlisting}

This is a simple formula for the fine-structure constant given the two
stiffness parameters $\kappa_{U(1)}$ and $\kappa_{SU(2)}$:
\begin{equation}
\alpha = \frac{e^2}{4\pi}, \qquad
\frac{1}{e^2} = \kappa_{U(1)} + \frac{\kappa_{SU(2)}}{4}.
\end{equation}

\textbf{This function takes $\kappa_{U(1)}$ as INPUT, not computing it.}

%=============================================================================
\section{The Existing \texttt{compute\_alpha.py} Script}
%=============================================================================

The file \texttt{compute\_alpha.py} (written by a previous agent) is a
415-line script that implements the full derivation chain. Its key
computational steps are:

\subsection{Step 1: Chern--Simons Weighted Average (Exact Finite Sum)}

The weight function and average:
\begin{equation}
w(k) = \frac{2}{k+2}\sin^2\!\frac{\pi}{k+2},
\qquad
\beta_{U(1)} = \frac{\sum_{k=0}^{59}(k+2)\,w(k)}{\sum_{k=0}^{59}w(k)}
= 3.796945276\ldots
\end{equation}

This is computed numerically as a finite sum of 60 terms. Each term involves
$\sin^2(\pi/(k+2))$, so this is NOT a simple rational expression. The sum
converges to a specific irrational number.

\subsection{Step 2: The Approximate Formula (85 ppm precision)}

The script computes an approximate closed-form:
\begin{equation}\label{eq:approx}
\boxed{
\alpha^{-1} \approx \frac{35\pi}{3} \times \beta_{U(1)}
  \times \frac{63}{64} \times \frac{4096}{4095}
}
\end{equation}
where the prefactor decomposes as:
\begin{equation}
\frac{35\pi}{3} = 4\pi \times \underbrace{\frac{5}{2}}_{\text{EW mixing}}
  \times \underbrace{\frac{7}{6}}_{\text{spectral triple}},
\qquad
\frac{5}{2} = 1 + \frac{R_W}{4}, \quad
\frac{7}{6} = \frac{N_{\rm species}}{N_{\rm gen}\cdot n_2}.
\end{equation}

Numerically:
\begin{align}
\alpha^{-1}_{\rm approx} &= 36.6519\ldots \times 3.7969\ldots
  \times 0.984375 \times 1.000244\ldots \\
&= 137.0244\ldots
\end{align}
This misses the experimental value by $\sim 85$~ppm.

\subsection{Step 3: The ``Exact'' Route (Reverse-Engineered)}

For the sub-ppm result, the script does NOT compute $\kappa_{U(1)}$
from first principles. Instead it computes:
\begin{equation}
\kappa_{U(1)} = \frac{\alpha^{-1}_{\rm target}}{6\pi}
= \frac{137.03599985}{6\pi} = 7.269986\ldots
\end{equation}

\textbf{This is circular}---it derives $\kappa$ FROM the target $\alpha^{-1}$,
not the other way around. The script explicitly labels this as
``Route A: From target $\alpha^{-1} = 137.03599985$ (reverse-engineering).''

The script also states (lines 220--234):
\begin{quote}
\textit{``The DFD papers state that the spectral action on the
Toeplitz-truncated $\mathbb{C}P^2 \times S^3$ geometry produces
$\kappa_{U(1)} = 7.26999\ldots$ through numerical evaluation of the
$a_4$ Seeley--DeWitt coefficient. This is NOT a simple closed-form
expression---it involves the full heat-kernel expansion.''}
\end{quote}

%=============================================================================
\section{The Missing Computation}
%=============================================================================

\subsection{What the Theory Claims}

The v108 Unified Theory (\texttt{section\_alpha\_locked.tex}) describes
the following locked derivation chain:
\begin{enumerate}
\item Operator: connection Laplacian $P = -g^{ij}\nabla_i\nabla_j$ on
  internal microsector $X = \mathbb{C}P^2 \times S^3$
\item Regulator: plateau cutoff (all negative moments $f_{-2} = f_{-4} = \cdots = 0$)
\item Finite-$k$ truncation: Toeplitz quantization at level $d = k+4 = 64$
\item Extract gauge kinetic term from $a_4$ Seeley--DeWitt coefficient
\item Microsector trace: regular module $\mathcal{H}_F = M_d(\mathbb{C})$,
  giving BOOST factor $\varepsilon = d^2/(d^2-1) = 4096/4095$
\end{enumerate}

The claim is that this chain produces $\kappa_{U(1)} = 7.26999\ldots$
and hence $\alpha^{-1} = 6\pi \times 7.26999\ldots = 137.03599985$.

\subsection{What Has Been Computed by Previous Agents}

Several Python scripts in the repository attempt the spectral action
computation:
\begin{itemize}
\item \texttt{clean\_a4\_computation.py} --- 500+ line script computing
  $a_4$ for the scalar Laplacian and Dirac operator on $\mathbb{C}P^2 \times S^3$
  via curvature invariants and the product convolution formula.
\item \texttt{numerical\_a4.py} --- Numerical extraction of $a_4$ from
  truncated eigenvalue sums with polynomial fitting.
\item \texttt{product\_a4.py} --- Explicit product formula
  $a_4(X) = \sum_{p+q=4} a_p(M_1)\,a_q(M_2)$.
\item \texttt{toeplitz\_heat\_kernel.py} --- Attempts to model Toeplitz
  truncation effects on the spectrum.
\end{itemize}

None of these scripts successfully reproduce $\kappa_{U(1)} = 7.26999\ldots$
from first principles. They compute the \emph{continuum} $a_4$ coefficient
correctly, but the critical step---the Toeplitz truncation at $d = 64$
and its effect on the spectral action---remains unresolved.

\subsection{The Kappa--Beta Gap}

The ratio $C = \kappa_{U(1)}/\beta_{U(1)} = 7.26999/3.79695 = 1.91469\ldots$
encodes all the spectral geometry. In the approximate formula
(Eq.~\ref{eq:approx}), this ratio is approximated by:
\begin{equation}
C_{\rm approx} = \frac{35}{18} \times \frac{63}{64} \times \frac{4096}{4095}
= 1.91453\ldots
\end{equation}
which differs from the exact value by $\sim 85$~ppm. The brute-force search
in \texttt{derive\_kappa.py} tried hundreds of rational multiples of $\pi$
and SM integers but found no exact closed-form match.

%=============================================================================
\section{Classification of the Computation}
%=============================================================================

\begin{center}
\renewcommand{\arraystretch}{1.4}
\begin{tabular}{lp{8cm}c}
\toprule
\textbf{Type} & \textbf{Description} & \textbf{Present?} \\
\midrule
(a) Lattice Monte Carlo &
  The Ab Initio paper describes lattice MC simulations at
  $(\beta_{U(1)}, \beta_{SU(2)}) = (3.80, 22.80)$ that yield
  $\alpha \approx 1/137$ within $\sim 1\%$. This is a
  \emph{verification}, not the sub-ppm derivation. & Verification only \\
(b) Spectral action / heat kernel &
  Claimed as the source of $\kappa_{U(1)} = 7.26999\ldots$ in
  the v108 theory, but the actual numerical evaluation of $a_4$
  on the Toeplitz-truncated geometry is never shown. & Claimed, not shown \\
(c) Direct algebraic formula &
  The approximate formula Eq.~(\ref{eq:approx}) achieves 85~ppm.
  For sub-ppm, the value $\kappa_{U(1)}$ is taken as a stated result. & Yes (approx.) \\
(d) Something else &
  The sub-ppm number $137.03599985$ originates from the combination
  $6\pi \times \kappa_{U(1)}$ where $\kappa$ is asserted from the
  spectral action. The actual computation is not published. & --- \\
\bottomrule
\end{tabular}
\end{center}

%=============================================================================
\section{The Exact Formulas That ARE Present}
%=============================================================================

\subsection{Formula 1: The Electroweak Relation (Exact)}

\begin{equation}
\boxed{\alpha^{-1} = 6\pi\,\kappa_{U(1)}}
\end{equation}

Derivation:
\begin{align}
\frac{1}{e^2} &= \kappa_{U(1)} + \frac{\kappa_{SU(2)}}{4}
  & \text{(electroweak mixing)} \\
\kappa_{SU(2)} &= 2\,\kappa_{U(1)}
  & \text{(Frame Stiffness Theorem: } \kappa_{U(1)}/\kappa_{SU(2)} = n_1/n_2 = 1/2\text{)} \\
\frac{1}{e^2} &= \frac{3}{2}\,\kappa_{U(1)} \\
\alpha^{-1} &= \frac{4\pi}{e^2} = 6\pi\,\kappa_{U(1)}
\end{align}

\subsection{Formula 2: The Chern--Simons Average (Exact Finite Sum)}

\begin{equation}
\boxed{
\beta_{U(1)} = \frac{\displaystyle\sum_{k=0}^{59}(k+2)\,w(k)}
                     {\displaystyle\sum_{k=0}^{59} w(k)},
\qquad
w(k) = \frac{2}{k+2}\sin^2\!\frac{\pi}{k+2}
}
\end{equation}

Numerically: $\beta_{U(1)} = 3.796945276\ldots$

\subsection{Formula 3: The Approximate Prefactor (85 ppm)}

\begin{equation}
\boxed{
\alpha^{-1} \approx \frac{35\pi}{3}\;\beta_{U(1)}
  \;\frac{63}{64}\;\frac{4096}{4095}
= 137.0244\ldots
}
\end{equation}

\subsection{Formula 4: The Lattice Estimator (1\% precision)}

From Listing~2 of the paper:
\begin{equation}
\alpha_W = \frac{g_1^2\,g_2^2}{(g_1^2 + g_2^2)\,4\pi},
\qquad g_1^2 = 1/\kappa_{U(1)}, \quad g_2^2 = 4/\kappa_{SU(2)}
\end{equation}

This is used in the lattice MC verification with
$\kappa_{U(1)} = \beta_{U(1)}$ and $\kappa_{SU(2)} = \beta_{SU(2)}$
measured from Wilson loops.

%=============================================================================
\section{Critical Assessment}
%=============================================================================

\subsection{What the Code Actually Shows}

\begin{enumerate}
\item The Chern--Simons weighted average $\beta_{U(1)} = 3.7969\ldots$
  is computed exactly from a well-defined 60-term finite sum. This is
  reproducible and unambiguous.

\item The UV cutoff $k_{\max} = 60$ is derived from the Bridge Lemma
  ($\chi(\mathbb{C}P^2, \mathcal{O}(9) \oplus \mathcal{O}^{\oplus 5}) = 60$),
  which is rigorous algebraic geometry.

\item The stiffness ratio $\kappa_{U(1)}/\kappa_{SU(2)} = 1/2$ and the
  Wilson ratio $\beta_{SU(2)}/\beta_{U(1)} = 6$ follow from the Frame
  Stiffness Theorem and $N_{\rm gen} = 3$.

\item The lattice MC verification at $(\beta_{U(1)}, \beta_{SU(2)}) = (3.80, 22.80)$
  yields $\alpha \approx 1/137$ within $\sim 1\%$. This is genuine
  numerical evidence but at $\sim 1\%$ precision, not sub-ppm.

\item \textbf{The sub-ppm value $\kappa_{U(1)} = 7.26999\ldots$ is an
  asserted output of a spectral action computation that is never shown
  in any code or publication.} The value $137.03599985$ follows trivially
  from $6\pi \times 7.26999\ldots$
\end{enumerate}

\subsection{The Gap}

The missing piece is the explicit computation:
\begin{equation}
\kappa_{U(1)} = \text{(spectral action on Toeplitz-truncated }
\mathbb{C}P^2 \times S^3\text{)} = 7.26999\ldots
\end{equation}

Until this computation is published with verifiable intermediate steps,
the sub-ppm claim rests on an unverified assertion. The approximate
formula achieves only 85~ppm agreement, which is far from the claimed
sub-ppm precision.

The approximate formula can be rewritten as:
\begin{equation}
\alpha^{-1}_{\rm approx} = 6\pi \times
\underbrace{\frac{35}{18}}_{\text{spectral geometry?}}
\times \beta_{U(1)} \times \frac{63}{64} \times \frac{4096}{4095}
\end{equation}
where $35/18 = 1.9444\ldots$ approximates the true ratio
$C = \kappa/(\beta \cdot (63/64) \cdot (4096/4095)) = 1.9147\ldots$
The 1.6\% discrepancy between $35/18$ and $C$ is what produces the
85~ppm residual.

\end{document}
